\scanpage{308}

\mspage{49}
% archive sheet number: 307

ce qui \deleted{signifie} revient à la tautologie suivante qu'un chemin \inserted{$C$} dans $X$ de type $\tau\tau'$ ``n'est autre chose'' qu'un couple $(c,c')$ d'un chemin \inserted{$c$} de type $\tau$ et d'un chemin $c'$ de type $\tau'$ ``mis bout à bout'' (dans cet ordre), i.e. tels que
\[
\operatorname{but}(c)=\operatorname{source}(c')
\]
(et on a alors
\[
\begin{aligned}
\operatorname{\unsure{source}}(C)&=\operatorname{source}(c)\\
\operatorname{but}(C)&=\operatorname{but}(c')\qquad ).
\end{aligned}
\]

Il se pose ici le sempiternel pb si on va noter ce chemin $c.c'$, ce qui semble à première vue s'imposer, ou plutôt $c'\circ c$ (\deleted{\unsure{en}} \unsure{s'imaginant|imaginant} \deleted{$c$}\inserted{le chemin $c$ de $x=\operatorname{source}(c)$ vers $y=\operatorname{but}(c)$} comme une \deleted{application} sorte de morphisme de $x$ \deleted{\unsure{vers}} $y$ \deleted{\unsure{et}} (cf plus bas le formalisme des cat de fractions), $c'$ comme un morphisme de $y=\deleted{\unsure{so}}\operatorname{source}(c')$ vers $z=\operatorname{but}(c')$, et en \unsure{interprétant} le composé $C$ comme le composé de ces deux morphismes \inserted{$c,c'$} (dans cet ordre du point de vue des opérations à effectuer), qu'on note traditionnellement $c'\circ c$ ou $c'.c$, et non $c\circ c'$ ou $c.c'$, hélas! C'est cette notation que nous allons suivre plus bas.

Mais pour l'instant on ne va pas se soucier des chemins dans des cat. quelconques $X$, mais revenir au fonctorialisme des $\tau$ et des $I_\tau$ eux-mêmes. L'interprétation catégorique de $\unsure{\tau\tau'|I_{\tau'}}$ comme une somme amalgamée nous assure que cette opération de composition est fonctorielle :
